\documentclass{beamer}

\usepackage[utf8]{inputenc}

\usepackage{color}

\usepackage{listings}

\usepackage{german}

\usecolortheme{beaver}
\definecolor{darkred}{RGB}{155,29,35}
\setbeamertemplate{navigation symbols}{}
\setbeamercolor{itemize item}{fg=darkred}
\setbeamercolor{itemize subitem}{fg=darkred}
\setbeamercolor{itemize subsubitem}{fg=darkred}

\definecolor{dkgreen}{rgb}{0,0.5,0}
\definecolor{gray}{rgb}{0.5,0.5,0.5}
\definecolor{mauve}{rgb}{0.58,0,0.82}
\definecolor{codebackground}{rgb}{0.9,0.9,0.9}

\lstdefinelanguage{scala}{
  morekeywords={abstract,case,catch,class,def,
    do,else,extends,false,final,finally,
    for,if,implicit,import,match,mixin,
    new,null,object,override,package,
    private,protected,requires,return,sealed,
    super,this,throw,trait,true,try,
    type,val,var,while,with,yield},
  otherkeywords={<-,<\%,<:,>:,\#,@},
  sensitive=true,
  morecomment=[l]{//},
  morecomment=[n]{/*}{*/},
  morestring=[b]",
  morestring=[b]',
  morestring=[b]"""
}

\lstdefinestyle{terminal}{
  language=,
  keywordstyle=,
  commentstyle=,
  stringstyle=
}

\lstset{
  language=scala,
  aboveskip=3mm,
  belowskip=3mm,
  showstringspaces=false,
  columns=flexible,
  basicstyle={\small\ttfamily},
  keywordstyle=\color{darkred},
  commentstyle=\color{gray},
  stringstyle=\color{dkgreen},
  numberstyle=\tiny\color{gray},
  numbers=none,
  numbersep=5pt,
  frame=none,
  breaklines=true,
  breakatwhitespace=true,
  tabsize=2,
  backgroundcolor=\color{codebackground}
}

\setlength\fboxsep{0pt}

\title{Ein Griff in die Scala-Trickkiste}
\author{Heiko Seeberger, Roman Roelofsen}
\institute{WeigleWilczek}
\date{JAX, Mainz, 2. Mai 2011}

\begin{document}


% Titlepage %%%%%%%%%%%%%%%%%%%%%%%%%%%%%%%%%%%%%%%%%%%%%%%%%%%%%%%%%%%%%%%%%%%%%%%%%%%%%%%%%%%%%%%%
\begin{frame}
  \titlepage
\end{frame}


% Agenda %%%%%%%%%%%%%%%%%%%%%%%%%%%%%%%%%%%%%%%%%%%%%%%%%%%%%%%%%%%%%%%%%%%%%%%%%%%%%%%%%%%%%%%%%%%
\begin{frame}
  \frametitle{Agenda}
  \tableofcontents
\end{frame}



% Section %%%%%%%%%%%%%%%%%%%%%%%%%%%%%%%%%%%%%%%%%%%%%%%%%%%%%%%%%%%%%%%%%%%%%%%%%%%%%%%%%%%%%%%%%%
\section{Builder Pattern}


\begin{frame}[fragile]{Builder Pattern}
\begin{itemize}
	\item API Design zum Setzen mehrere Parameter
	\item setter-Methode-Ansatz
	\begin{lstlisting}
		val config = new Config
		config.setA(1)
		config.enableLog()
		config.setB(2)
		val created = config.create()
	\end{lstlisting}
	\item Builder-Pattern-Ansatz
	\begin{lstlisting}
		val created =
		  new Config().setA(1).setB(2).enableLog().create()
	\end{lstlisting}	
\end{itemize}
\end{frame}

%%%%%%%%%%%%%%%%%%%%%%%%%%%%%%%%%%%%%%%%%%%%%%%%%%%%%%%%%%%%%%%%%%%%%%%%%%%%%%%%%%

\begin{frame}[fragile]{Builder Pattern}
\begin{itemize}
	\item Methodenaufrufe geben i.d.R. \textbf{this} zurück...
	\item ... oder ein Objekt welches sich die Konfiguration "`merkt"'
	\item Manche Konfiguration stehen eventuell erst nach dem Setzen bestimmer Werte zur Verfügung
	\begin{lstlisting}
		val created = new Config().enableLog().withDbLog().create()
	\end{lstlisting}
\end{itemize}
\end{frame}  
  
%%%%%%%%%%%%%%%%%%%%%%%%%%%%%%%%%%%%%%%%%%%%%%%%%%%%%%%%%%%%%%%%%%%%%%%%%%%%%%%%%%

\begin{frame}[fragile]{Builder Pattern - Implementierung}
\begin{itemize}
	\item Definition
		\begin{lstlisting}
		class ABad {
		  var number = 1
  
		  def double(): ABad = {
		    number *= 2
		    this
		  }
		  def triple(): ABad = {
		    number = number * 3
		    this
		  }
		}
		\end{lstlisting}
	\item Verwendung
		\begin{lstlisting}
		println(new ABad().double().triple().number)
		\end{lstlisting}
\end{itemize}
\end{frame}

%%%%%%%%%%%%%%%%%%%%%%%%%%%%%%%%%%%%%%%%%%%%%%%%%%%%%%%%%%%%%%%%%%%%%%%%%%%%%%%%%%

\begin{frame}[fragile]{Builder Pattern - Implementierung}
\begin{itemize}
	\item Definition
		\begin{lstlisting}
		class BBad1 extends ABad {
		  def random(): BBad1 = {
		    number *= Random.nextInt
		    this
		  }
		}
		\end{lstlisting}
	\item Verwendung
		\begin{lstlisting}
		println(new BBad1().random().number)
		println(new BBad1().random().triple().number)
		\end{lstlisting}
	\item Problem
		\begin{lstlisting}
		println(new BBad1().double().triple().random().number)
		\end{lstlisting}
\end{itemize}
\end{frame}

%%%%%%%%%%%%%%%%%%%%%%%%%%%%%%%%%%%%%%%%%%%%%%%%%%%%%%%%%%%%%%%%%%%%%%%%%%%%%%%%%%

\begin{frame}[fragile]{Builder Pattern - Implementierung}
\begin{itemize}
	\item Lösung (wie in Java)
		\begin{lstlisting}
		class BBad2 extends ABad {
		  override def double(): BBad2 = {
		    super.double()
		    this
		  }

		  override def triple(): BBad2 = {
		    super.triple()
		    this
		  }

		  def random(): BBad2 = {
		    number *= Random.nextInt
		    this
		  }
		}
		\end{lstlisting}
\end{itemize}
\end{frame}

%%%%%%%%%%%%%%%%%%%%%%%%%%%%%%%%%%%%%%%%%%%%%%%%%%%%%%%%%%%%%%%%%%%%%%%%%%%%%%%%%%

\begin{frame}[fragile]{Builder Pattern - Implementierung}
\begin{itemize}
	\item Lösung (wie in Java)
		\begin{itemize}
			\item Alle Methodendeklarationen müssen wiederholt werden
			\item Lösung ist sehr fragil, da Änderungen in der Superklasse (z.B. neue Methode) die Kette "`unterbrechen"' können
			\item Alle Subklassen müssen informiert und nachgezogen werden
		\end{itemize}
\end{itemize}
\end{frame}

%%%%%%%%%%%%%%%%%%%%%%%%%%%%%%%%%%%%%%%%%%%%%%%%%%%%%%%%%%%%%%%%%%%%%%%%%%%%%%%%%%

\begin{frame}[fragile]{Builder Pattern - Implementierung}
\begin{itemize}
	\item Lösung mit Scala this.type
		\begin{lstlisting}
		class AGood {
		  var number = 1
		  def double(): this.type = {
		    number *= 2
		    this
		  }
		  def triple(): this.type = {
		    number = number * 3
		    this
		  }
		}
		\end{lstlisting}
\end{itemize}
\end{frame}

%%%%%%%%%%%%%%%%%%%%%%%%%%%%%%%%%%%%%%%%%%%%%%%%%%%%%%%%%%%%%%%%%%%%%%%%%%%%%%%%%%

\begin{frame}[fragile]{Builder Pattern - Implementierung}
\begin{itemize}
	\item Lösung mit Scala this.type
		\begin{lstlisting}
		class BGood extends AGood {
		  def random(): this.type = {
		    number *= Random.nextInt
		    this
		  }
		}
		\end{lstlisting}
	\item Verwendung
		\begin{lstlisting}
		println(new BGood().double().triple().random().number)
		\end{lstlisting}
\end{itemize}
\end{frame}


% Section %%%%%%%%%%%%%%%%%%%%%%%%%%%%%%%%%%%%%%%%%%%%%%%%%%%%%%%%%%%%%%%%%%%%%%%%%%%%%%%%%%%%%%%%%%
\section{Stackable Modifications}

\begin{frame}[fragile]{Stackable Modifications}
\begin{itemize}
	\item Logging Klassen
		\begin{itemize}
			\item Basis-Interface: \verb!Log!
			\item Implementierungen: \verb!NormalLog!, \verb!DBLog!
		\end{itemize}
		\pause
	\item Modifizierungen
		\begin{itemize}
			\item \verb!ImportantLog! (Text + '!!!')
			\item \verb!SuperImportantLog! (Großbuchstaben)
		\end{itemize}
		\pause
	\item Ziel
		\begin{itemize}
			\item Log-Implementierungen sollen beliebig mit den Modifizierungen erweitert werden können
			\item Modifizierungen müssen sich kombinieren lassen
			\item Je nach Komplexität kann die Reihenfolge wichtig sein
		\end{itemize}
\end{itemize}
\end{frame}

%%%%%%%%%%%%%%%%%%%%%%%%%%%%%%%%%%%%%%%%%%%%%%%%%%%%%%%%%%%%%%%%%%%%%%%%%%%%%%%%%%

\begin{frame}[fragile]{Stackable Modifications - Log Klassen}
\begin{lstlisting}
abstract class Log {
  def log(msg: String): Unit
}

class ConsoleLog extends Log {
  def log(msg: String) {
    println("LOG: " + msg)
  }
}

class DBLog extends Log {
  def log(msg: String) {
    Database.save(msg)
  }
}
\end{lstlisting}
\end{frame}

%%%%%%%%%%%%%%%%%%%%%%%%%%%%%%%%%%%%%%%%%%%%%%%%%%%%%%%%%%%%%%%%%%%%%%%%%%%%%%%%%%

\begin{frame}[fragile]{Stackable Modifications - Modifizierungen}
\begin{lstlisting}
trait ImportantLog extends Log {
  abstract override def log(msg: String) {
    super.log(msg +  "!!!")
  }
}

trait SuperImportantLog extends Log {
  abstract override def log(msg: String) {
    super.log(msg.toUpperCase)
  }
}
\end{lstlisting}
\end{frame}

%%%%%%%%%%%%%%%%%%%%%%%%%%%%%%%%%%%%%%%%%%%%%%%%%%%%%%%%%%%%%%%%%%%%%%%%%%%%%%%%%%

\begin{frame}[fragile]{Stackable Modifications - Verwendung}
\begin{lstlisting}
class AppLog extends ConsoleLog with ImportantLog with SuperImportantLog
val log1 = new AppLog
log1.log("A log message")

val log2 = new ConsoleLog with ImportantLog
log2.log("Another message")

val log3 = new DBLog with SuperImportantLog
log3.log("Another message")

val log4 = new DBLog with ImportantLog with SuperImportantLog
log4.log("Another message")
\end{lstlisting}
\end{frame}


% Section %%%%%%%%%%%%%%%%%%%%%%%%%%%%%%%%%%%%%%%%%%%%%%%%%%%%%%%%%%%%%%%%%%%%%%%%%%%%%%%%%%%%%%%%%%
\section{Kontrollstrukturen abstrahieren}


\begin{frame}[fragile]{Kontrollstrukturen abstrahieren}
\center \textbf{Aufgabe} \linebreak	
(1) Resource öffnen (2) Resource verwenden \linebreak (3) Resource schließen (4) Fehler abfangen
\center \textbf{Java}
\begin{lstlisting}[language=java]
InputStream is = null;
try {
  is = new FileInputStream(new File("..."));
  is.read();
} catch (Exception e) {
  // ...
} finally {
  try {
    is.close();
  } catch (Exception e1) {
    // ...
  }
}
\end{lstlisting}
\end{frame}  

%%%%%%%%%%%%%%%%%%%%%%%%%%%%%%%%%%%%%%%%%%%%%%%%%%%%%%%%%%%%%%%%%%%%%%%%%%%%%%%%%%

\begin{frame}[fragile]{Kontrollstrukturen abstrahieren}
\center \textbf{Aufgabe} \linebreak	
(1) Resource öffnen (2) Resource verwenden \linebreak (3) Resource schließen (4) Fehler abfangen
\center \textbf{Scala}
\begin{lstlisting}
withResource(new FileInputStream(new File("..."))) { is =>
  // ...
}
\end{lstlisting}
\end{frame}  

%%%%%%%%%%%%%%%%%%%%%%%%%%%%%%%%%%%%%%%%%%%%%%%%%%%%%%%%%%%%%%%%%%%%%%%%%%%%%%%%%%

\begin{frame}[fragile]{Kontrollstrukturen abstrahieren}
\center \textbf{Aufgabe} \linebreak	
(1) Resource öffnen (2) Resource verwenden \linebreak (3) Resource schließen (4) Fehler abfangen
\center \textbf{Scala Implementierung}
\begin{lstlisting}
def withResource[A <: { def close() }](res: A)(block: A => Any) {
  try {
    block(res)
  } catch {
    case e: Throwable => throw e
  } finally {
    try {
      res.close()
    } catch {
      case e: Throwable => throw e
    }
  }
}
\end{lstlisting}
\end{frame}  

%%%%%%%%%%%%%%%%%%%%%%%%%%%%%%%%%%%%%%%%%%%%%%%%%%%%%%%%%%%%%%%%%%%%%%%%%%%%%%%%%%

\begin{frame}[fragile]{Kontrollstrukturen abstrahieren}
\begin{itemize}
	\item Java-Code verwendet oft Exceptions als gültigen Rückgabewert
	\item Automatische Exception / Option Konvertierung im Scala-Kontext daher sinnvoll
\end{itemize}
\begin{lstlisting}
import scala.util.control.Exception._

val saveParse = catching(classOf[NumberFormatException])
val numberOption = saveParse opt { "123b".toInt }   

numberOption match {
  case Some(n) => n
  case None => 0
}
\end{lstlisting}
\end{frame}  


% Section %%%%%%%%%%%%%%%%%%%%%%%%%%%%%%%%%%%%%%%%%%%%%%%%%%%%%%%%%%%%%%%%%%%%%%%%%%%%%%%%%%%%%%%%%%
\section{Implicits und Type Classes}

\begin{frame}
  \frametitle{Implicits und Type Classes}
  \includegraphics[width=\linewidth]{img/free-climbing.jpg}
\end{frame}

\begin{frame}[fragile]
  \frametitle{Ein sauberes Design}
    \begin{lstlisting}
object Grade {
  object Qualifier extends Enumeration {
    val Plus = Value("+")
    val Minus = Value("-")
  }
}

case class Grade(
    value: Int, 
    qualifier: Option[Grade.Qualifier.Value] = None)

case class Route(name: String, grade: Grade)
    \end{lstlisting}
\end{frame}

\begin{frame}[fragile]
  \frametitle{Wo liegt das Problem?}
  \begin{itemize}
    \item M\"uhsam zu benutzen!
    \begin{lstlisting}
println(Route("Fight Gravity", Grade(8, Some(Plus))))
println(Route("Kasperltheater", Grade(8)))
    \end{lstlisting}
    \item W\"are das nicht sch\"oner?
    \begin{lstlisting}
println(Route("Fight Gravity", Grade(8, Plus)))
println(Route("Kasperltheater", 8))
    \end{lstlisting}
    \item Aber das compiliert doch nicht, oder?
  \end{itemize}
\end{frame}

\begin{frame}[fragile]
  \frametitle{Zum erwarteten Typ konvertieren}
  \begin{itemize}
    \item Zum Gl\"uck gibt der Scala-Compiler nicht gleich auf!
    \item Der Ausweg hei{\ss}t \emph{Implicit Conversion to Expected Type}:
    \begin{lstlisting}
object Grade {

  object Qualifier extends Enumeration {
    ...
    implicit def valueToOption(value: Value) =
        Option(value)
  }

  implicit def intToGrade(value: Int) = Grade(value)
}
    \end{lstlisting}
  \end{itemize}
\end{frame}

\begin{frame}[fragile]
  \frametitle{Was k\"onnten wir sonst noch besser machen?}
  \begin{itemize}
    \item Annahme: Wir sollen unsere Objekte nach XML serialisieren
    \item Schlechte Idee: Klassen mit \emph{toXml}-Methoden "`verschmutzen"'
    \item Besser: Wir rufen einfach \emph{toXml} auf unseren Objekten auf!
    \begin{lstlisting}
println(Route("Kasperltheater", 8).toXml)
    \end{lstlisting}
    \item Aber das compiliert doch schon wieder nicht, oder?
  \end{itemize}
\end{frame}

\begin{frame}[fragile]
  \frametitle{Pimp my Library}
  \begin{itemize}
    \item Der Scala-Compiler gibt auch hier noch nicht auf!
    \item Gesucht wird: \emph{Implicit Conversion of the Receiver}
    \item Oder: Wer macht mir aus einer \emph{Route} etwas mit \emph{toXml}?
    \begin{lstlisting}
object Route {
  implicit def toXmlSerializable(route: Route) = new {
    def toXml =
      <route name={ route.name } >{ 
        route.grade.toXml
      }</route>
  }
}
    \end{lstlisting}
  \end{itemize}
\end{frame}

\begin{frame}[fragile]
  \frametitle{Aber geht es nicht noch besser?}
  \begin{itemize}
    \item W\"are es nicht noch sch\"oner, wenn es \emph{toXml} f\"ur "`alles"' g\"abe?
    \item Sodass wir ... 
    \item ... zumindest die Signatur abstrahieren k\"onnen
    \item ... und flexibel bei der Implementierung bleiben
  \end{itemize}
\end{frame}

\begin{frame}[fragile]
  \frametitle{Type Classes machen es m\"oglich!}
  \begin{itemize}
    \item Erst einmal abstrahieren wir \emph{toXml}:
    \begin{lstlisting}
trait XmlSerializable[A] {
  implicit def toXmlSerializable(a: A) = new {
    def toXml(implicit format: XmlFormat[A]) =
      format toXml a
  }
}

trait XmlFormat[A] {
  def toXml(a: A): NodeSeq
}
    \end{lstlisting}
    \item Wir delegieren also zum impliziten \emph{XmlFormat[A]}
    \item Aber woher bekommen wir "`Instanzen"' f\"ur diese \emph{Type Class}
  \end{itemize}
\end{frame}

\begin{frame}[fragile]
  \frametitle{Type Classes implementieren}
  \begin{itemize}
    \item Dazu definieren wir implizite Werte/Objekte:
    \begin{lstlisting}
object Route extends XmlSerializable[Route] {
  implicit object RouteXmlFormat extends
      XmlFormat[Route] {
    def toXml(route: Route) =
      <route name={ route.name } >{
        route.grade.toXml }
      </route>
  }
}
    \end{lstlisting}
  \end{itemize}
\end{frame}


% Section %%%%%%%%%%%%%%%%%%%%%%%%%%%%%%%%%%%%%%%%%%%%%%%%%%%%%%%%%%%%%%%%%%%%%%%%%%%%%%%%%%%%%%%%%%
\section{Delimited Continuations}

\begin{frame}
  \frametitle{Delimited Continuations}
  \includegraphics[width=\linewidth]{img/continue.jpg}
\end{frame}

\begin{frame}[fragile]
  \frametitle{Continuation Passing Style}
  \begin{itemize}
    \item Gew\"ohnlich ruft die "`sp\"atere"' Funktion die "`fr\"uhere"' auf:
    \begin{lstlisting}
def plusOne(i: Int) = i + 1
def minusOne(i: Int) = i - 1
println(minusOne(plusOne(0)))
    \end{lstlisting}
    \item Aber es geht auch "`anders herum"':
    \begin{lstlisting}
def plusOneCps(i: Int, cont: Int => Int) = cont(i + 1)
println(plusOneCps(0, minusOne))
    \end{lstlisting}
    \item "`Sp\"atere"' Funktion als Argument der "`fr\"uheren"'!
  \end{itemize}
\end{frame}

\begin{frame}[fragile]
  \frametitle{Invertierung des Kontrollflusses}
  \begin{itemize}
    \item Die "`fr\"uhere"' Funktion entscheidet, wie es weiter geht:
    \begin{lstlisting}
def plusOneCps2(i: Int, cont: Int => Int) =
  if (i > 0) cont(i + 1) else Int.MinValue
    \end{lstlisting}
    \item Das kann sehr n\"utzlich sein, f\"uhrt jedoch zu schwer verst\"andlichem Code f\"ur den Aufrufer:
    \begin{lstlisting}
println(plusOneCps2(0, minusOne))
    \end{lstlisting}
  \end{itemize}
\end{frame}

\begin{frame}[fragile]
  \frametitle{Scala macht es wieder einfach!}
  \begin{itemize}
    \item Mit \emph{shift} greifen wir uns eine Continuation:
    \begin{lstlisting}
def plusOneCont2(i: Int) =
  shift { cont: (Int => Int) =>
    if (i > 0) cont(i + 1) else Int.MinValue
  }
    \end{lstlisting}
    \item Mit \emph{reset} geben wir die Grenzen vor:
    \begin{lstlisting}
println(reset { minusOne(plusOneCont2(0)) })
    \end{lstlisting}
  \end{itemize}
\end{frame}

\begin{frame}[fragile]
  \frametitle{Beispiel: Eingabe-Validierung I}
  \begin{itemize}
    \item Eigentlich haben wir einen einfachen Kontrollfluss:
    \begin{lstlisting}
val x = read("First Int: ")
val y = read("Second Int: ")
val result = x + y
println("Result: %s" format result)

def read(label: String) = readLine(label).toInt
    \end{lstlisting}
    \item Aber: Eingabe-Fehler m\"ussten vom Aufrufer behandelt werden
    \item Der Kontrollfluss wird kompliziert!
  \end{itemize}
\end{frame}

\begin{frame}[fragile]
  \frametitle{Beispiel: Eingabe-Validierung II}
  \begin{itemize}
    \item Verlagern der Validierung in die \emph{read}-Methode:
    \begin{lstlisting}
def read(label: String) =
  shift { cont: (Int => Unit) =>
    val s = readLine(label)
    try {
      cont(s.toInt)
    } catch {
      case _: NumberFormatException =>
        println("Illegal argument: %s" format s)
    }
  }
    \end{lstlisting}
    \item Dadurch bleibt der Kontrollfluss einfach!
  \end{itemize}
\end{frame}



% End %%%%%%%%%%%%%%%%%%%%%%%%%%%%%%%%%%%%%%%%%%%%%%%%%%%%%%%%%%%%%%%%%%%%%%%%%%%%%%%%%%%%%%%%%%

\begin{frame}
  \frametitle{Danke!}
  \begin{itemize}
    \item Folien: \href{https://github.com/downloads/weiglewilczek/advanced-scala/advanced-scala.pdf}{github.com/downloads/weiglewilczek/advanced-scala/advanced-scala.pdf}
    \item Code: \href{https://github.com/weiglewilczek/advanced-scala}{github.com/weiglewilczek/advanced-scala}
    \item Twitter: \href{http://twitter.com/hseeberger}{hseeberger} und \href{http://twitter.com/romanroe}{romanroe}
  \end{itemize}
\end{frame}


\end{document}
